% 1. ročník MAT – příprava 99; inspirace Realisticky.cz 020501–020506, úlohy nově sestavené.
\documentclass[11pt,a4paper]{article}
\input{../../sablony/pripravy/priprava.sty}
\usepackage{tikz}
\usepackage{pgfplots}
\pgfplotsset{compat=1.18}
\usetikzlibrary{arrows.meta}
\nastavPripravu{1. ročník}{Kvadratické funkce}{99}{Doplnění na čtverec}
\begin{document}
\blok{Proč potřebujeme doplnění na čtverec?}
U vrcholového tvaru poznáme vrchol okamžitě. Ale předpis $f(x)=x^2-6x+5$ jej neukazuje. Vrátíme se k algebraickému vzorci $(u-v)^2=u^2-2uv+v^2$. Přičteme a odečteme stejné číslo, takže se funkce nezmění.
\begin{definiceBox}[title={Definice - Doplnění na čtverec}]
\textbf{Doplnění na čtverec} je ekvivalentní úprava
\[ x^2+px+q=\left(x+\frac p2\right)^2+q-\left(\frac p2\right)^2. \]
Pro $a\ne0$ nejdříve vytkneme $a$ z členů obsahujících $x$; cílem je předpis ve vrcholovém tvaru.
\end{definiceBox}
\blok{Společně řešené příklady}

\begin{prikladReseniBox}[title={Příklad 1}]
Upravte $f(x)=x^2-6x+5$ na vrcholový tvar a určete vrchol, minimum a obor hodnot.
\tcblower
\[
\begin{aligned}
f(x)&=x^2-6x+9-9+5 \\
&=(x-3)^2-4 \\
V&=[3;-4],\quad \min f=-4 \\
H(f)&=\langle-4;+\infty)
\end{aligned}
\]

\end{prikladReseniBox}
\begin{prikladReseniBox}[title={Příklad 2}]
Doplňte na čtverec $g(x)=x^2+4x+1$. Proč musíme přidat i odečíst právě číslo $4$?
\tcblower
\[
\begin{aligned}
g(x)&=x^2+4x+4-4+1 \\
&=(x+2)^2-3 \\
V&=[-2;-3],\quad H(g)=\langle-3;+\infty)
\end{aligned}
\]
Platí $4=2^2$, protože prostřední člen $4x=2\cdot x\cdot2$.
\end{prikladReseniBox}
\begin{prikladReseniBox}[title={Příklad 3}]
Upravte $h(x)=2x^2-8x+3$. Nezapomeňte vytknout dvojku ještě před doplněním čtverce.
\tcblower
\[
\begin{aligned}
h(x)&=2(x^2-4x)+3 \\
&=2\bigl((x-2)^2-4\bigr)+3 \\
&=2(x-2)^2-5 \\
V&=[2;-5],\quad H(h)=\langle-5;+\infty)
\end{aligned}
\]

\end{prikladReseniBox}
\blok{Společný krátký postup s negativním $a$}
U $k(x)=-x^2+6x-7$ nejprve vytkneme mínus:
\[\begin{aligned}k(x)&=-(x^2-6x)-7\\&=-\bigl((x-3)^2-9\bigr)-7\\&=-(x-3)^2+2.\end{aligned}\]
Vrchol je $[3;2]$, hodnota $2$ je \emph{maximum}.
\blok{Samostatné procvičování}

\begin{zadaniReseniBox}[title={Samostatná práce 1}]
a) $x^2-4x+1$; b) $x^2+2x-3$; c) $-x^2-2x+3$. Upravte na vrcholový tvar a zapište vrchol.
\tcblower
\[
\begin{aligned}
x^2-4x+1&=(x-2)^2-3\quad V[2;-3] \\
x^2+2x-3&=(x+1)^2-4\quad V[-1;-4] \\
-x^2-2x+3&=-(x+1)^2+4\quad V[-1;4]
\end{aligned}
\]

\end{zadaniReseniBox}
\blok{Domácí úkol}
\begin{zadaniReseniBox}[title={Domácí úkol}]
\begin{enumerate}[label=\alph*),itemsep=3pt]
\item Doplňte na čtverec $x^2-2x+3$ a určete vrchol.
\item Doplňte na čtverec $x^2+6x+5$ a určete minimum.
\item Upravte $x^2-8x+10$ a určete $H(f)$.
\item Upravte $-x^2+4x+1$ na vrcholový tvar a určete maximum.
\item Upravte $2x^2+8x+5$ na vrcholový tvar.
\item Upravte $-2x^2+8x-3$ na vrcholový tvar.
\item \textbf{Výzva:} Doplňte na čtverec $x^2+3x+1$ a určete vrchol.
\end{enumerate}
\tcblower
\begin{enumerate}[label=\alph*),itemsep=4pt]
\item
\[
\begin{aligned}
x^2-2x+3&=(x-1)^2+2 \\
V&=[1;2]
\end{aligned}
\]

\item
\[
\begin{aligned}
x^2+6x+5&=(x+3)^2-4 \\
\min f&=-4
\end{aligned}
\]

\item
\[
\begin{aligned}
x^2-8x+10&=(x-4)^2-6 \\
H(f)&=\langle-6;+\infty)
\end{aligned}
\]

\item
\[
\begin{aligned}
-x^2+4x+1&=-(x-2)^2+5 \\
\max f&=5
\end{aligned}
\]

\item
\[
\begin{aligned}
2x^2+8x+5&=2(x^2+4x)+5 \\
&=2(x+2)^2-3
\end{aligned}
\]

\item
\[
\begin{aligned}
-2x^2+8x-3&=-2(x^2-4x)-3 \\
&=-2(x-2)^2+5
\end{aligned}
\]

\item
\[
\begin{aligned}
x^2+3x+1&=(x+\tfrac32)^2-\tfrac94+1 \\
&=(x+\tfrac32)^2-\tfrac54 \\
V&=[-\tfrac32;-\tfrac54]
\end{aligned}
\]

\end{enumerate}
\end{zadaniReseniBox}
\poznamka{Čas: zopakování vzorce 5 min, dva monické příklady 12 min, koeficient před $x^2$ 10 min, případ se záporným $a$ 5 min, samostatně 10 min, závěr 3 min. V této hodině počítej bez vzorce pro kořeny kvadratické rovnice.}
\end{document}
